% Integrated 4 September 2026 from the isolated exact-parametrization audit.
% Every label is namespaced koide-param-*; the original variables, chamber,
% strict walls, lattice, and height are unchanged.
\providecommand{\OO}{\mathcal O}
\providecommand{\NN}{\mathbb N}
\providecommand{\ZZ}{\mathbb Z}
\providecommand{\RR}{\mathbb R}
\providecommand{\QQ}{\mathbb Q}
\providecommand{\Norm}{\operatorname N}
\providecommand{\Tr}{\operatorname {Tr}}
\providecommand{\e}{\varepsilon}

\section{Exact ideal and unit-orbit parametrization of the unchanged sharp Koide count}

Throughout this section the variables, equation, chamber, and height are
exactly
\begin{equation}
 a^2-3q^2-18n^2=-2,
 \qquad
 a>0,\quad q\geq0,\quad n\geq0,\quad a+4n>q,
 \qquad
 4a+17n\leq X.
 \label{koide-param-eq:original-problem}
\end{equation}
Every auxiliary symbol introduced below is connected to these same integer
points by an explicitly proved map and inverse.  No count for a normalized
form, a rescaled lattice, or a substituted cutoff is used in place of
\eqref{koide-param-eq:original-problem}.
Thus
\begin{equation}
 R(X)=\#\left\{(a,q,n)\in\ZZ^3:\text{all conditions in
 \eqref{koide-param-eq:original-problem} hold}\right\}.
 \label{koide-param-eq:R-definition}
\end{equation}
The already calibrated leading constant is
\begin{equation}
 C_0=
 \frac9{2\pi^2}\log(2+\sqrt3)
 \log\!\left(\frac3{2\sqrt2}\right).
 \label{koide-param-eq:C0}
\end{equation}

The exact parametrization proved below has four pieces of data for each
fixed $n\geq1$: a positive integer $C$, a positive integer $D$, a square
root $E$ of $3$ modulo $D$, and a nonnegative Pell exponent $k$.  There is
no quotient left implicit: equations \eqref{koide-param-eq:bounded-seed-search},
\eqref{koide-param-eq:HNF-to-seed}, and \eqref{koide-param-eq:unit-orbit-coordinates} construct
$(a,q,n)$ by finite integer operations, while
\eqref{koide-param-eq:forward-HNF} and the descent \eqref{koide-param-eq:inverse-descent} recover
$(C,D,E,k)$ from $(a,q,n)$.  The face $n=0$ is recorded separately and
exactly.

This parametrization gives the exact sharp-cutoff formula
\eqref{koide-param-eq:exact-floor-formula}.  It also separates the discrepancy into
the radial shifted-divisor sum and the unit-phase discrepancy in
\eqref{koide-param-eq:exact-discrepancy-splitting}.  The square-root estimate
\begin{equation}
 R(X)-C_0X=O_\e(X^{1/2+\e})
 \label{koide-param-eq:square-root-target}
\end{equation}
is not proved or refuted here.  What is proved unconditionally, by applying
the effective symmetric-variety theorem to the exact chamber, is that
there is a constant $\eta>0$ for which
\begin{equation}
 R(X)=C_0X+O(X^{1-\eta}).
 \label{koide-param-eq:unconditional-power-saving}
\end{equation}
The cited general theorem does not supply $\eta\geq1/2$, and its proof does
not justify replacing \eqref{koide-param-eq:unconditional-power-saving} by
\eqref{koide-param-eq:square-root-target}.

\section{Arithmetic of $\ZZ[\sqrt3]$}

Put
\begin{equation}
 K=\QQ(\sqrt3),\qquad \OO=\ZZ[\sqrt3],\qquad
 \epsilon=2+\sqrt3.
 \label{koide-param-eq:field-data}
\end{equation}
For $\alpha=u+v\sqrt3\in K$, write
\begin{equation}
 \alpha'=u-v\sqrt3,
 \qquad
 \Norm(\alpha)=\alpha\alpha'=u^2-3v^2.
 \label{koide-param-eq:conjugate-norm}
\end{equation}

\begin{lemma}[Ring, class number, and totally positive units]
\label{koide-param-lem:ring-units}
The ring of integers of $K$ is $\OO$.  Every nonzero integral ideal of
$\OO$ is principal.  Every unit has norm $+1$, and the totally positive
unit group is
\begin{equation}
 \OO^{\times,+}=\{\epsilon^j:j\in\ZZ\}.
 \label{koide-param-eq:positive-units}
\end{equation}
\end{lemma}

\begin{proof}
Since $3\equiv3\pmod4$, the standard integral-basis calculation for a
quadratic field gives $\OO_K=\ZZ[\sqrt3]$ and discriminant $12$.
Minkowski's bound says that every ideal class in a real quadratic field of
discriminant $12$ has an integral representative of norm at most
\(\frac{2!}{2^2}\sqrt{12}=\sqrt3<2\).  Its norm must therefore be $1$,
so it is $\OO$ and the ordinary class number is one.

The norm of a unit is $\pm1$.  The equation $u^2-3v^2=-1$ is impossible
modulo $3$, because it would give $u^2\equiv2\pmod3$.  Hence every unit
has norm $+1$.  Dirichlet's unit theorem makes the positive units an
infinite cyclic group.  To identify its least generator, suppose
$\gamma=u+v\sqrt3>1$ and $\gamma'>0$.  Replacing $\gamma$ by its inverse
if necessary makes $v\geq0$.  If $v=0$, then $\gamma=1$.  If $v\geq1$,
then $u^2-3v^2=1$ forces $u\geq2$, and consequently
$\gamma\geq2+\sqrt3=\epsilon$.  Since $\epsilon$ itself has norm one,
it is the least positive unit greater than one and generates the group.
\end{proof}

For $m\geq1$ define
\begin{equation}
 \chi_{12}(d)=\left(\frac{12}{d}\right),
 \qquad
 \rho(m)=\sum_{d\mid m}\chi_{12}(d).
 \label{koide-param-eq:rho-definition}
\end{equation}
Here the symbol is Kronecker's symbol; in particular
\begin{equation}
 \chi_{12}(d)=
 \begin{cases}
  0,&2\mid d\text{ or }3\mid d,\\
  1,&d\equiv1,11\pmod {12},\\
  -1,&d\equiv5,7\pmod {12}.
 \end{cases}
 \label{koide-param-eq:chi12-table}
\end{equation}

\begin{lemma}[Ideal number]
\label{koide-param-lem:ideal-number}
The number of integral ideals of $\OO$ having absolute norm $m$ is
$\rho(m)$.
\end{lemma}

\begin{proof}
The Dedekind zeta function factors as
$\zeta_K(s)=\zeta(s)L(s,\chi_{12})$.  Comparing Dirichlet coefficients
gives the assertion.  Equivalently, the equality follows prime by prime.
At $2$ and $3$ there is one ramified prime ideal, so there is one ideal of
norm $p^e$ for every $e\geq0$.  If $\chi_{12}(p)=1$, then $p$ splits and
there are $e+1$ ideals of norm $p^e$.  If $\chi_{12}(p)=-1$, then $p$ is
inert, so there is one such ideal when $e$ is even and none when $e$ is
odd.  These are respectively the local values of
$\sum_{j=0}^e\chi_{12}(p)^j$.
Consequently, if
$m=2^{e_2}3^{e_3}\prod_{p\geq5}p^{e_p}$, then
\begin{equation}
 \rho(m)=
 \begin{cases}
  \displaystyle\prod_{\substack{p\mid m,\ p\geq5\\p\equiv1,11\ (12)}}(e_p+1),
   &e_p\text{ is even for every }p\equiv5,7\pmod {12},\\[6pt]
  0,&\text{otherwise}.
 \end{cases}
 \label{koide-param-eq:rho-Euler-product}
\end{equation}
\end{proof}

\section{Hermite-normal-form parameters for every ideal}

For $m\geq1$ let
\begin{equation}
 \mathcal P(m)=
 \left\{(C,D,E):
 \begin{array}{l}
 C,D\in\ZZ_{>0},\quad C^2D=m,\\
 0\leq E<D,\quad E^2\equiv3\pmod D
 \end{array}
 \right\}.
 \label{koide-param-eq:P-m}
\end{equation}

\begin{proposition}[Exact HNF parametrization of ideals]
\label{koide-param-prop:HNF-ideals}
The map
\begin{equation}
 (C,D,E)\longmapsto
 \mathfrak a(C,D,E)
 :=C\bigl(D\ZZ+(E+\sqrt3)\ZZ\bigr)
 \label{koide-param-eq:param-ideal}
\end{equation}
is a bijection from $\mathcal P(m)$ to the integral ideals of $\OO$ of
norm $m$.  Consequently
\begin{equation}
 \#\mathcal P(m)
 =\sum_{C^2D=m}\#\{E\bmod D:E^2\equiv3\pmod D\}
 =\rho(m).
 \label{koide-param-eq:root-divisor-identity}
\end{equation}
\end{proposition}

\begin{proof}
Relative to the basis $1,\sqrt3$ of $\OO$, the lattice in
\eqref{koide-param-eq:param-ideal} has column basis
\[
 \begin{pmatrix}CD&CE\\0&C\end{pmatrix}.
\]
Its index is $C^2D=m$.  It is stable under multiplication by $\sqrt3$
because
\begin{align*}
 \sqrt3(CD)&=D\,C(E+\sqrt3)-E(CD),\\
 \sqrt3\,C(E+\sqrt3)
 &=E\,C(E+\sqrt3)+\frac{3-E^2}{D}(CD),
\end{align*}
and the last coefficient is integral.  It is therefore an integral ideal
of norm $m$.

Conversely, write the unique column Hermite normal form of an integral
ideal $\mathfrak a$ as
\[
 \mathfrak a=A\ZZ+(B+C\sqrt3)\ZZ,
 \qquad A,C>0,\quad0\leq B<A.
\]
Stability of $A$ under multiplication by $\sqrt3$ implies $C\mid A$ and
$C\mid B$.  Write $A=CD$ and $B=CE$.  Stability of $B+C\sqrt3$ then
implies $D\mid3-E^2$.  The HNF inequality gives $0\leq E<D$, and the
index is $AC=C^2D$.  HNF uniqueness proves injectivity and surjectivity.

For completeness, the last equality in \eqref{koide-param-eq:root-divisor-identity}
can be checked without zeta functions.  Let
$r_3(D)=\#\{E\bmod D:E^2\equiv3\pmod D\}$.  For $p\geq5$, Hensel's
lemma gives $r_3(p^j)=2$ for every $j\geq1$ if
$\chi_{12}(p)=1$, and $r_3(p^j)=0$ if $\chi_{12}(p)=-1$.  Also
$r_3(2)=r_3(3)=1$ and $r_3(2^j)=r_3(3^j)=0$ for $j\geq2$.
Thus, for every prime power $p^e$,
\[
 \sum_{2c+j=e}r_3(p^j)=\sum_{j=0}^e\chi_{12}(p)^j.
\]
Multiplicativity proves \eqref{koide-param-eq:root-divisor-identity}.
\end{proof}

The HNF parameters can be recovered directly from a generator.  Suppose
$\alpha=a+q\sqrt3$ has positive norm $m$, with $a>0$ and $q\geq0$.  Put
\begin{equation}
 C=\gcd(a,q),\qquad A=a/C,\qquad B=q/C,\qquad D=m/C^2.
 \label{koide-param-eq:forward-HNF-pre}
\end{equation}
If $D=1$, put $E=0$.  If $D>1$, then
\begin{equation}
 \gcd(B,D)=1,\qquad E\equiv A B^{-1}\pmod D,\qquad0\leq E<D.
 \label{koide-param-eq:forward-HNF}
\end{equation}
Indeed,
$\gcd(B,D)=\gcd(B,A^2-3B^2)=\gcd(B,A^2)=1$, and
$E^2\equiv3\pmod D$.  The principal ideal $(\alpha)$ is exactly
$\mathfrak a(C,D,E)$: the element $A+B\sqrt3$ lies in
$D\ZZ+(E+\sqrt3)\ZZ$, so the principal ideal is contained in the HNF
ideal, while both ideals have norm $D$ after division by $C$.

\section{Complete integral parametrization of the original equation}

For $n\geq1$ retain the exact integer
\begin{equation}
 m_n=18n^2-2.
 \label{koide-param-eq:mn}
\end{equation}
It satisfies $m_n>0$ and $m_n\equiv1\pmod3$.

\begin{lemma}[Every ideal of norm $m_n$ has positive-norm generators]
\label{koide-param-lem:positive-norm-generator}
If $\mathfrak a$ is an integral ideal of $\OO$ of norm $m_n$, then every
generator has norm $+m_n$, rather than $-m_n$.
\end{lemma}

\begin{proof}
Lemma~\ref{koide-param-lem:ring-units} makes $\mathfrak a$ principal.  A generator
$u+v\sqrt3$ has norm $\pm m_n$.  Negative norm would give
\[
 u^2-3v^2=-m_n,
\]
and reduction modulo $3$ would give
$u^2\equiv-m_n\equiv2\pmod3$, which is impossible.  Thus the norm is
$+m_n$.
\end{proof}

For each $\mathbf p=(C,D,E)\in\mathcal P(m_n)$, choose a generator of
$\mathfrak a(C,D,E)$.  By Lemma~\ref{koide-param-lem:positive-norm-generator} its
norm is $m_n$.  Change its sign so that both real embeddings are positive,
and multiply by a unique power of $\epsilon$ so as to obtain the unique
generator
\begin{equation}
 \alpha_{n,\mathbf p}=a_{n,\mathbf p,0}
   +q_{n,\mathbf p,0}\sqrt3,
 \qquad
 \sqrt{m_n}\leq\alpha_{n,\mathbf p}
       <\epsilon\sqrt{m_n}.
 \label{koide-param-eq:canonical-generator}
\end{equation}
Existence and uniqueness follow from
\eqref{koide-param-eq:positive-units}, because the half-open intervals
$[\epsilon^j\sqrt{m_n},\epsilon^{j+1}\sqrt{m_n})$ partition
$\RR_{>0}$.  Since
$\alpha_{n,\mathbf p}'=m_n/\alpha_{n,\mathbf p}$, the lower inequality
in \eqref{koide-param-eq:canonical-generator} implies
$q_{n,\mathbf p,0}\geq0$.

\begin{lemma}[Purely integral reduced seed and descent]
\label{koide-param-lem:integral-seed-descent}
For an integral pair $(a,q)$ satisfying
\[
 a^2-3q^2=m_n,\qquad a>0,\qquad q\geq0,
\]
the half-open inequalities
\[
 \sqrt{m_n}\leq a+q\sqrt3<\epsilon\sqrt{m_n}
\]
are equivalent to the purely integral inequalities
\begin{equation}
 a>2q\geq0.
 \label{koide-param-eq:integral-seed-condition}
\end{equation}
If $a\leq2q$, multiplication by $\epsilon^{-1}=2-\sqrt3$ is the
integral inverse step
\begin{equation}
 (a,q)\longmapsto(a^-,q^-)=(2a-3q,2q-a).
 \label{koide-param-eq:inverse-descent}
\end{equation}
It preserves the norm, has $a^->0$ and $0\leq q^-<q$, and can therefore
be iterated until the unique pair satisfying \eqref{koide-param-eq:integral-seed-condition}
is reached.
\end{lemma}

\begin{proof}
Put $u=(a+q\sqrt3)/\sqrt{m_n}$.  Norm positivity and the sign conditions
give $u\geq1$.  Moreover
\[
 \frac a{\sqrt{m_n}}=\frac{u+u^{-1}}2,
 \qquad
 \frac{q\sqrt3}{\sqrt{m_n}}=\frac{u-u^{-1}}2.
\]
For $u\geq1$, the ratio
\[
 \frac aq=\sqrt3\,\frac{u^2+1}{u^2-1}
\]
(with value $+\infty$ when $q=0$) is strictly decreasing.  At
$u=\epsilon$ it equals $2$, since
$\epsilon+\epsilon^{-1}=4$.  Thus $u<\epsilon$ is equivalent to $a>2q$;
the strict inequality exactly matches the half-open upper endpoint.

If $a\leq2q$, norm positivity gives $a>\sqrt3q$, and hence
\[
 a^-=2a-3q>(2\sqrt3-3)q>0,
 \qquad q^-=2q-a\geq0.
\]
Also $a>\sqrt3q>q$ gives $q^-=2q-a<q$.  Thus the positive integer $q$
strictly decreases at every inverse step, so the process terminates.  A
terminal pair has $a>2q$.  Conversely, for such a pair the second
coordinate after one inverse step is negative, so no further inverse step
stays in the nonnegative-$q$ region.  This proves uniqueness as well as
termination.
\end{proof}

The preceding seed can be reconstructed from the HNF data by a bounded
integer search; thus the parametrization does not require an unspecified
choice of a generator.  More precisely, for
$\mathbf p=(C,D,E)\in\mathcal P(m_n)$ there is a unique pair $(A,B)$ with
\begin{equation}
 \begin{gathered}
  A,B\in\ZZ,\qquad A>0,\qquad 0\leq B,\qquad B^2<D,\\
  A^2-3B^2=D,\qquad A\equiv EB\pmod D.
 \end{gathered}
 \label{koide-param-eq:bounded-seed-search}
\end{equation}
It gives the generator in \eqref{koide-param-eq:canonical-generator} by the exact
integer formula
\begin{equation}
 (a_{n,\mathbf p,0},q_{n,\mathbf p,0})=(CA,CB).
 \label{koide-param-eq:HNF-to-seed}
\end{equation}
Indeed, existence follows by writing the reduced generator as
$C(A+B\sqrt3)$.  The seed inequality $A>2B$ gives
$D=A^2-3B^2>B^2$.  Conversely, the conditions in
\eqref{koide-param-eq:bounded-seed-search} imply $A>2B$ and put
$A+B\sqrt3$ in $D\ZZ+(E+\sqrt3)\ZZ$.  Its principal ideal is contained
in that HNF ideal, and both have norm $D$, so they are equal.  Lemma
\ref{koide-param-lem:integral-seed-descent} and HNF uniqueness give uniqueness.  In
particular one needs to test only the integers
$0\leq B\leq\lfloor\sqrt{D-1}\rfloor$ and the positive square
$A=\sqrt{D+3B^2}$.
The phase introduced below is therefore also explicit in the same integer
data:
\begin{equation}
 \tau_{n,(C,D,E)}
 =\frac{\log((A+B\sqrt3)/\sqrt D)}{\log\epsilon}.
 \label{koide-param-eq:HNF-phase-explicit}
\end{equation}
Here the factor $C$ cancels because $m_n=C^2D$; no ideal generator remains
to be chosen.

For $k\geq0$ define integers $a_{n,\mathbf p,k}$ and
$q_{n,\mathbf p,k}$ by
\begin{equation}
 a_{n,\mathbf p,k}+q_{n,\mathbf p,k}\sqrt3
 =\alpha_{n,\mathbf p}\epsilon^k.
 \label{koide-param-eq:unit-orbit-coordinates}
\end{equation}
Equivalently,
\begin{equation}
 \begin{pmatrix}a_{n,\mathbf p,k+1}\\q_{n,\mathbf p,k+1}\end{pmatrix}
 =
 \begin{pmatrix}2&3\\1&2\end{pmatrix}
 \begin{pmatrix}a_{n,\mathbf p,k}\\q_{n,\mathbf p,k}\end{pmatrix}.
 \label{koide-param-eq:unit-recurrence}
\end{equation}

\begin{theorem}[Integral parametrization, with inverse and multiplicity]
\label{koide-param-thm:integral-parametrization}
The map
\begin{equation}
 (n,\mathbf p,k)
 \longmapsto
 (a_{n,\mathbf p,k},q_{n,\mathbf p,k},n),
 \qquad
 n\geq1,\quad\mathbf p\in\mathcal P(m_n),\quad k\geq0,
 \label{koide-param-eq:parameter-map}
\end{equation}
is a bijection onto all integral solutions of
\begin{equation}
 a^2-3q^2-18n^2=-2,\qquad a>0,\ q\geq0,\ n\geq1.
 \label{koide-param-eq:n-positive-solutions}
\end{equation}
Every such solution automatically satisfies $a+4n>q$.

The inverse is as follows.  From $(a,q,n)$ form
$\alpha=a+q\sqrt3$ and recover $(C,D,E)$ by
\eqref{koide-param-eq:forward-HNF-pre}--\eqref{koide-param-eq:forward-HNF}.  There is a unique
$k\geq0$ such that
\begin{equation}
 \alpha=\alpha_{n,(C,D,E)}\epsilon^k.
 \label{koide-param-eq:recover-k}
\end{equation}
Thus every point occurs with multiplicity exactly one.
\end{theorem}

\begin{proof}
Equation \eqref{koide-param-eq:unit-orbit-coordinates} preserves the norm, so its
image satisfies
\[
 a_{n,\mathbf p,k}^2-3q_{n,\mathbf p,k}^2=m_n,
\]
which is exactly the equation in \eqref{koide-param-eq:n-positive-solutions}.
Both embeddings of the generator are positive, hence $a>0$; and
$\alpha_{n,\mathbf p}\geq\sqrt{m_n}$ and $k\geq0$ imply $q\geq0$.

Conversely, a solution in \eqref{koide-param-eq:n-positive-solutions} has
\(\Norm(a+q\sqrt3)=m_n>0\).  Since $a>0$ and $q\geq0$, the inequality
$a^2-3q^2=m_n>0$ gives $a>\sqrt3q$; hence both embeddings are positive.
Its principal ideal has unique HNF data $(C,D,E)$.  The quotient of
 $a+q\sqrt3$ by the canonical generator of that ideal is a totally positive
 unit, hence is $\epsilon^k$ for a unique $k\in\ZZ$.  Lemma
 \ref{koide-param-lem:integral-seed-descent} supplies the completely integral recovery
 of $k$: iterate \eqref{koide-param-eq:inverse-descent} precisely while $a\leq2q$;
 the number of steps is $k$, and the terminal pair is
 $(a_{n,\mathbf p,0},q_{n,\mathbf p,0})$.  In particular $k\geq0$.  This
 proves surjectivity and the stated inverse.  HNF uniqueness, seed
 uniqueness, and unit-exponent uniqueness prove multiplicity one.

It remains to prove the chamber wall.  If a point with $n\geq1$ had
$q\geq a+4n$, then $a\leq q-4n$, and the unchanged equation would give
\[
 3q^2+18n^2-2=a^2\leq(q-4n)^2.
\]
After moving terms and dividing by $2$, this is
\[
 q^2+4qn+n^2\leq1.
\]
For nonnegative integers with $n\geq1$, equality could occur only at
$(q,n)=(0,1)$, but that pair does not satisfy $q\geq a+4n$ with $a>0$.
Thus the assumed wall inequality is impossible.
\end{proof}

The congruence needed by the original inverse map to the $(x,y,z)$
coordinates is retained automatically: reducing
$a^2-3q^2-18n^2=-2$ modulo $2$ gives
\begin{equation}
 a\equiv q\pmod2.
 \label{koide-param-eq:parity}
\end{equation}
No parity class has been lost or added in
Theorem~\ref{koide-param-thm:integral-parametrization}.

\begin{proposition}[The face $n=0$]
\label{koide-param-prop:n-zero}
All points satisfying the original equation and chamber with $n=0$ are
\begin{equation}
 a_k+q_k\sqrt3=(1+\sqrt3)\epsilon^k,\qquad k\geq1.
 \label{koide-param-eq:n-zero-family}
\end{equation}
They have height $4a_k$ and occur with multiplicity one.
\end{proposition}

\begin{proof}
When $n=0$, the norm is $-2$.  Its parity implies $a\equiv q\pmod2$, and
therefore
\[
 \frac{a+q\sqrt3}{1+\sqrt3}
 =\frac{3q-a}{2}+\frac{a-q}{2}\sqrt3\in\OO.
\]
The quotient has norm one, so it is $\pm\epsilon^k$ by
Lemma~\ref{koide-param-lem:ring-units}.  The conditions $a>0$ and $q\geq0$ select the
positive sign.  At $k=0$ one gets $(a,q)=(1,1)$, which is excluded by the
strict chamber wall $a>q$.  At every $k\geq1$, the recurrence
\eqref{koide-param-eq:unit-recurrence} gives
$a_{k+1}-q_{k+1}=a_k+q_k>0$.  For $k\leq-1$, direct use of
$a=(t-2/t)/2$ with $t=(1+\sqrt3)\epsilon^k\leq\sqrt3-1<\sqrt2$
gives $a<0$.  Hence precisely the exponents $k\geq1$ occur.
\end{proof}

\section{Exact divisor and sharp-height formulas}

The number of unit-orbits at a fixed positive $n$ has an especially simple
divisor formula.

\begin{proposition}[Orbit divisor sum and shifted factorization]
\label{koide-param-prop:orbit-divisor}
For $n\geq1$, the number of orbits under
\((a,q)\mapsto(2a+3q,a+2q)\) among the solutions with $a>0$ and
$q\geq0$ is
\begin{equation}
 \rho(m_n)=\sum_{d\mid18n^2-2}\chi_{12}(d)
 =\rho(3n-1)\rho(3n+1).
 \label{koide-param-eq:orbit-divisor-sum}
\end{equation}
\end{proposition}

\begin{proof}
The first equality is Theorem~\ref{koide-param-thm:integral-parametrization} and
Proposition~\ref{koide-param-prop:HNF-ideals}.  For the second, write
\[
 m_n=2(3n-1)(3n+1).
\]
The two displayed linear factors have greatest common divisor dividing
$2$.  The function $\rho=1*\chi_{12}$ is multiplicative and
$\rho(2^j)=1$ for all $j\geq0$.  Removing all powers of $2$ therefore
makes the two factors coprime without changing any of the three $\rho$
values, proving the product identity.
\end{proof}

For $\mathbf p\in\mathcal P(m_n)$ define its exact unit phase by
\begin{equation}
 \tau_{n,\mathbf p}
 =\frac{\log(\alpha_{n,\mathbf p}/\sqrt{m_n})}{\log\epsilon},
 \qquad0\leq\tau_{n,\mathbf p}<1.
 \label{koide-param-eq:unit-phase}
\end{equation}
For $n\geq1$ put
\begin{equation}
 Y_n(X)=\frac{X-17n}{4}.
 \label{koide-param-eq:Y-n}
\end{equation}
Whenever $Y_n(X)\geq\sqrt{m_n}$, put
\begin{equation}
 L_n(X)=
 \frac1{\log\epsilon}
 \log\!\left(
 \frac{Y_n(X)+\sqrt{Y_n(X)^2-m_n}}{\sqrt{m_n}}
 \right).
 \label{koide-param-eq:L-n}
\end{equation}
Let
\begin{equation}
 \mathcal N(X)=
 \{n\in\ZZ_{\geq1}:Y_n(X)\geq\sqrt{m_n}\}.
 \label{koide-param-eq:N-X}
\end{equation}
This is the exact finite range
\begin{equation}
 X\geq17n+4\sqrt{18n^2-2};
 \label{koide-param-eq:n-boundary}
\end{equation}
no asymptotic replacement of this boundary is made.

For the $n=0$ face define
\begin{equation}
 R_0(X)=
 \max\!\left\{0,
 \left\lfloor
 \frac{
 \log\!\left(
 \dfrac{X/4+\sqrt{(X/4)^2+2}}{1+\sqrt3}
 \right)}{\log\epsilon}
 \right\rfloor\right\}.
 \label{koide-param-eq:R0-exact}
\end{equation}

\begin{theorem}[Exact sharp-height floor formula]
\label{koide-param-thm:exact-floor}
For every real $X\geq0$,
\begin{equation}
 R(X)=R_0(X)+
 \sum_{n\in\mathcal N(X)}
 \sum_{\mathbf p\in\mathcal P(m_n)}
 \left(\left\lfloor L_n(X)-\tau_{n,\mathbf p}\right\rfloor+1\right).
 \label{koide-param-eq:exact-floor-formula}
\end{equation}
Every equality case in $4a+17n\leq X$ is included by the floor in
\eqref{koide-param-eq:exact-floor-formula}.
\end{theorem}

\begin{proof}
For a fixed positive $n$ and parameter $\mathbf p$, put
$t_k=\alpha_{n,\mathbf p}\epsilon^k$.  Since
$t_k\geq\sqrt{m_n}$,
\begin{equation}
 a_{n,\mathbf p,k}=\frac12\left(t_k+\frac{m_n}{t_k}\right)
 \label{koide-param-eq:a-from-t}
\end{equation}
is nondecreasing, and is strictly increasing after the possible initial
case $q=0$.  The unchanged height inequality is
$a_{n,\mathbf p,k}\leq Y_n(X)$.  It has no solution unless
$Y_n(X)\geq\sqrt{m_n}$.  In that range, solving the quadratic inequality
in $t_k$ gives
\begin{equation}
 t_k\leq Y_n(X)+\sqrt{Y_n(X)^2-m_n}.
 \label{koide-param-eq:t-height-bound}
\end{equation}
After division by $\sqrt{m_n}$ and logarithms, this is precisely
\begin{equation}
 k+\tau_{n,\mathbf p}\leq L_n(X).
 \label{koide-param-eq:k-phase-bound}
\end{equation}
Because $L_n(X)\geq0$ and $0\leq\tau_{n,\mathbf p}<1$, the number of
integers $k\geq0$ satisfying \eqref{koide-param-eq:k-phase-bound} is exactly
$\lfloor L_n(X)-\tau_{n,\mathbf p}\rfloor+1$; this expression is zero,
not negative, when $L_n(X)<\tau_{n,\mathbf p}$.  Equality is retained.
Theorem~\ref{koide-param-thm:integral-parametrization} now proves the positive-$n$
sum.

For $n=0$, put $t=(1+\sqrt3)\epsilon^k$.  Since the norm is $-2$,
$a_k=(t-2/t)/2$.  The inequality $4a_k\leq X$ is equivalent to
\[
 t\leq X/4+\sqrt{(X/4)^2+2}.
\]
Counting the integers $k\geq1$ gives \eqref{koide-param-eq:R0-exact} and completes
the proof.
\end{proof}

Write
\begin{equation}
 J_n(X)=\lfloor L_n(X)\rfloor,
 \qquad
 \lambda_n(X)=L_n(X)-J_n(X)\in[0,1),
 \label{koide-param-eq:J-lambda}
\end{equation}
and define the exact phase distribution
\begin{equation}
 A_n(t)=\#\{\mathbf p\in\mathcal P(m_n):
                     \tau_{n,\mathbf p}\leq t\},
 \qquad0\leq t<1,
 \label{koide-param-eq:A-n}
\end{equation}
with discrepancy
\begin{equation}
 \Delta_n(t)=A_n(t)-t\rho(m_n).
 \label{koide-param-eq:Delta-n}
\end{equation}
Since
\(\lfloor J+\lambda-\tau\rfloor+1
=J+\mathbf1_{\{\tau\leq\lambda\}}\), Theorem~\ref{koide-param-thm:exact-floor}
gives the exact identity
\begin{align}
 R(X)-C_0X
 &=R_0(X)
 +\left[
   \sum_{n\in\mathcal N(X)}
     \rho(m_n)L_n(X)-C_0X
  \right]
 +\sum_{n\in\mathcal N(X)}
    \Delta_n(\lambda_n(X)).
 \label{koide-param-eq:exact-discrepancy-splitting}
\end{align}
This identity includes the two unbounded coordinate faces and every sharp
height boundary point.  Formula \eqref{koide-param-eq:R0-exact} also gives
$R_0(X)=O(\log(2+X))$, so the $n=0$ face is below the requested
$X^{1/2+\varepsilon}$ scale; the two positive-$n$ brackets are the genuine
analytic terms.

\section{The phase term as exact Hecke-character data}

The last sum in \eqref{koide-param-eq:exact-discrepancy-splitting} is not determined
by the divisor count $\rho(m_n)$.  It records where the generator of each
ideal lies inside a fundamental interval for the positive units.

For $\ell\in\ZZ$ and a nonzero integral ideal $\mathfrak a=(\gamma)$,
define
\begin{equation}
 \xi_\ell(\mathfrak a)=
 \exp\!\left(
  \frac{\pi i\ell}{\log\epsilon}
  \log\left|\frac{\gamma}{\gamma'}\right|
 \right).
 \label{koide-param-eq:Hecke-character}
\end{equation}
This is independent of the generator: replacing $\gamma$ by
$\pm\epsilon^j\gamma$ changes the logarithm by $2j\log\epsilon$, and
the exponential by $e^{2\pi i\ell j}=1$.  It is multiplicative, hence is
the unramified ideal part of a unitary Hecke character of $K$.  For
$m=m_n$ and the canonical generator,
\begin{equation}
 \xi_\ell(\mathfrak a(C,D,E))
 =e^{2\pi i\ell\tau_{n,(C,D,E)}}.
 \label{koide-param-eq:character-phase}
\end{equation}
Define its norm-$m$ coefficient by
\begin{equation}
 S_\ell(m)=\sum_{\Norm\mathfrak a=m}\xi_\ell(\mathfrak a).
 \label{koide-param-eq:S-ell}
\end{equation}
Then $S_0(m)=\rho(m)$, whereas the nonzero $S_\ell(m_n)$ are the Fourier
coefficients of the phase multiset.

The coefficient $S_\ell(m)$ also has a completely explicit Euler-factor
description.  It is multiplicative in $m$.  If $p\geq5$ is split, choose
one prime ideal $\mathfrak p$ above $p$ and put
$\omega_{\ell,p}=\xi_\ell(\mathfrak p)$.  Since
$\mathfrak p\overline{\mathfrak p}=(p)$ and
$\xi_\ell((p))=1$, one has
\begin{equation}
 S_\ell(p^e)=\sum_{j=0}^{e}\omega_{\ell,p}^{2j-e}
 \qquad(\chi_{12}(p)=1).
 \label{koide-param-eq:split-Hecke-factor}
\end{equation}
If $p\geq5$ is inert, then
\begin{equation}
 S_\ell(p^e)=
 \begin{cases}1,&2\mid e,\\0,&2\nmid e.
 \end{cases}
 \qquad(\chi_{12}(p)=-1).
 \label{koide-param-eq:inert-Hecke-factor}
\end{equation}
Finally the unique primes above $2$ and $3$ are generated by
$1+\sqrt3$ and $\sqrt3$, respectively.  Since
\[
 \left|\frac{1+\sqrt3}{1-\sqrt3}\right|=\epsilon,
 \qquad
 \left|\frac{\sqrt3}{-\sqrt3}\right|=1,
\]
their factors are
\begin{equation}
 S_\ell(2^e)=(-1)^{\ell e},
 \qquad S_\ell(3^e)=1.
 \label{koide-param-eq:ramified-Hecke-factors}
\end{equation}
These formulas enumerate every ideal factor with multiplicity one.  At
$\ell=0$ they reduce prime by prime to the divisor sum $\rho(m)$; at
$\ell\ne0$ they show exactly what angular information the unweighted
divisor sum omits.

They also retain the shifted factorization of $m_n$ at every Fourier mode.
Writing $3n-1=2^r u$ and $3n+1=2^s v$ with $u,v$ odd, one has
$\gcd(u,v)=1$.  Multiplicativity and
\eqref{koide-param-eq:ramified-Hecke-factors} therefore give the exact identity
\begin{equation}
 S_\ell(18n^2-2)
 =(-1)^\ell S_\ell(3n-1)S_\ell(3n+1).
 \label{koide-param-eq:shifted-Hecke-factorization}
\end{equation}
Indeed, the exponent of $2$ on the left is $1+r+s$, while the product on
the right has total exponent $r+s$; all odd prime factors are disjoint.
Thus the angular obstruction is itself a shifted correlation, now of
nontrivial Hecke coefficients rather than merely of their dimensions.

Here is an endpoint-exact Fourier statement.  Suppose $0<\lambda<1$, and
let
\begin{equation}
 b_n(t)=\#\{\mathbf p\in\mathcal P(m_n):\tau_{n,\mathbf p}=t\}.
 \label{koide-param-eq:b-n}
\end{equation}
Fej\'er summation of the Fourier series of
$\mathbf1_{[0,\lambda]}-\lambda$ gives
\begin{align}
 \Delta_n(\lambda)
 &=\frac12\bigl(b_n(0)+b_n(\lambda)\bigr)\notag\\
 &\quad+
 \lim_{M\to\infty}
 \sum_{0<|\ell|\leq M}
 \left(1-\frac{|\ell|}{M+1}\right)
 \frac{1-e^{-2\pi i\ell\lambda}}{2\pi i\ell}
 S_\ell(m_n).
 \label{koide-param-eq:Fejer-exact}
\end{align}
The two half-weights correct the midpoint convention at the circular
endpoint $0$ and at $\lambda$.  When $\lambda=0$, the exact value is
simply $\Delta_n(0)=b_n(0)$.

Combining this with \eqref{koide-param-eq:shifted-Hecke-factorization} makes the
aggregate obstruction a literal shifted correlation.  Define the exact
endpoint contribution
\begin{equation}
 \mathcal E(X)=
 \sum_{\substack{n\in\mathcal N(X)\\\lambda_n(X)=0}}b_n(0)
 +\frac12
 \sum_{\substack{n\in\mathcal N(X)\\\lambda_n(X)>0}}
       \bigl(b_n(0)+b_n(\lambda_n(X))\bigr).
 \label{koide-param-eq:endpoint-aggregate}
\end{equation}
Since $\mathcal N(X)$ is finite, summation over $n$ commutes with the
Fej\'er limit and gives
\begin{align}
 \sum_{n\in\mathcal N(X)}\Delta_n(\lambda_n(X))
 &=\mathcal E(X)\notag\\
 &\quad+\lim_{M\to\infty}
 \sum_{0<|\ell|\leq M}
 \left(1-\frac{|\ell|}{M+1}\right)
 \frac{(-1)^\ell}{2\pi i\ell}\notag\\
 &\qquad\times
 \sum_{\substack{n\in\mathcal N(X)\\\lambda_n(X)>0}}
 \bigl(1-e^{-2\pi i\ell\lambda_n(X)}\bigr)
 S_\ell(3n-1)S_\ell(3n+1).
 \label{koide-param-eq:exact-shifted-Hecke-correlation}
\end{align}
Every coefficient and endpoint in this identity comes from the original
weak height inequality; no smoothing has been inserted.

Equations \eqref{koide-param-eq:exact-discrepancy-splitting} and
\eqref{koide-param-eq:Fejer-exact} identify the analytic problem without discarding
the original height.  Its radial part is the sharply weighted shifted
correlation
\begin{equation}
 \sum_{n\in\mathcal N(X)}
 \rho(3n-1)\rho(3n+1)L_n(X),
 \label{koide-param-eq:radial-shifted-correlation}
\end{equation}
and its angular part requires cancellation in the polynomially sampled
Hecke coefficients
\begin{equation}
 S_\ell(18n^2-2),
 \label{koide-param-eq:polynomial-Hecke-samples}
\end{equation}
uniformly over the Fourier modes needed to approximate the discontinuous
threshold $\tau\leq\lambda_n(X)$.  The endpoint
\eqref{koide-param-eq:n-boundary} gives $L_n(X)=0$ and has a square-root singularity
in its derivative; it is not a fixed smooth weight.  A bound for the
unweighted orbit total $\rho(m_n)$ therefore does not bound the actual
sharp-height count.

\section{An unconditional effective exponent for the exact sector}

This section records the power saving that is presently justified without
assuming the square-root estimate.

Let
\begin{equation}
 F(a,q,n)=a^2-3q^2-18n^2,
 \qquad
 \mathscr X=\{F=-2\},
 \qquad
 \mathbf G=\operatorname{Spin}(F).
 \label{koide-param-eq:group-variety}
\end{equation}
Let $x_0=(1,1,0)$.  Its orthogonal complement is
\begin{equation}
 x_0^\perp=\{(3s,s,t):s,t\in\QQ\},
 \qquad
 F(3s,s,t)=6(s^2-3t^2),
 \label{koide-param-eq:stabilizer-complement}
\end{equation}
so the identity component of its stabilizer is the norm-one torus of
$\QQ(\sqrt3)$.  Its rational character group is zero because Galois
conjugation acts by $-1$ on its geometric character lattice.

The form $F$ is anisotropic over $\QQ$.  Indeed, suppose a primitive
integral triple satisfies $F(a,q,n)=0$.  Modulo $3$ one first gets
$3\mid a$.  Writing $a=3a_1$ and dividing the equation by $3$ gives
$3a_1^2-q^2-6n^2=0$, hence $3\mid q$.  Write $q=3q_1$ and divide by
$3$ again to get
\begin{equation}
 a_1^2-3q_1^2-2n^2=0.
\end{equation}
Modulo $3$, the sum $a_1^2+n^2$ is zero, which forces
$3\mid a_1$ and $3\mid n$.  Thus $3$ divides $a,q,n$, contradicting
primitivity.  In particular the associated arithmetic Fuchsian quotient
is cocompact.

For $X>0$ let $B_X$ be the exact subset of $\mathscr X(\RR)$ cut out by
\begin{equation}
 a>0,\quad q\geq0,\quad n\geq0,\quad a+4n>q,
 \quad4a+17n\leq X.
 \label{koide-param-eq:B-X}
\end{equation}
With the gauge measure
\begin{equation}
 dm_\infty=\frac{dq\,dn}{2a}\qquad(a>0),
 \label{koide-param-eq:gauge-measure}
\end{equation}
put
\begin{equation}
 I_\infty=\frac1{2\sqrt3}
 \log\!\left(\frac3{2\sqrt2}\right).
 \label{koide-param-eq:I-infty}
\end{equation}

\begin{lemma}[Exact-sector volume with a power-saving remainder]
\label{koide-param-lem:volume}
As $X\to\infty$,
\begin{equation}
 m_\infty(B_X)=I_\infty X+O(1).
 \label{koide-param-eq:volume-linear-O1}
\end{equation}
Moreover, the family $B_X$ is effectively well-rounded in the sense of
Benoist--Oh, Definition~13.1~\cite[p.~1931]{BenoistOh_2012}.
\end{lemma}

\begin{proof}
Introduce exact polar coordinates on the negative plane by
\begin{equation}
 \sqrt3q=r\cos\theta,\qquad
 3\sqrt2n=r\sin\theta,\qquad
 0\leq\theta\leq\pi/2.
 \label{koide-param-eq:polar}
\end{equation}
On $F=-2$ one has $r=\sqrt{a^2+2}$ and
\begin{equation}
 dm_\infty=\frac1{6\sqrt6}\,da\,d\theta.
 \label{koide-param-eq:polar-measure}
\end{equation}
Write $b(\theta)=17\sin\theta/(3\sqrt2)$.  In the positive orthant the
outer boundary is the unique nonnegative solution $A_X(\theta)$ of
\begin{equation}
 4A_X(\theta)+b(\theta)\sqrt{A_X(\theta)^2+2}=X.
 \label{koide-param-eq:A-boundary}
\end{equation}
For all sufficiently large $X$, uniformly in $\theta$,
\begin{equation}
 A_X(\theta)=\frac{X}{4+b(\theta)}+O(X^{-1}).
 \label{koide-param-eq:A-asymptotic}
\end{equation}
To see this without a formal expansion, subtract
$(4+b)A_X$ from \eqref{koide-param-eq:A-boundary} and use
\[
 \sqrt{A_X^2+2}-A_X
 =\frac2{\sqrt{A_X^2+2}+A_X}=O(X^{-1}),
\]
where \eqref{koide-param-eq:A-boundary} uniformly gives $A_X\asymp X$.
Integration of \eqref{koide-param-eq:A-asymptotic} in
\eqref{koide-param-eq:polar-measure} gives
\[
 \frac{X}{6\sqrt6}\int_0^{\pi/2}
 \frac{d\theta}{4+\frac{17}{3\sqrt2}\sin\theta}+O(X^{-1})
 =I_\infty X+O(X^{-1}).
\]
For the asserted exact value, set
$c=17/(3\sqrt2)$, $d=\sqrt{c^2-16}=1/(3\sqrt2)$, and
$t=\tan(\theta/2)$.  Then
\begin{align*}
 \int_0^{\pi/2}\frac{d\theta}{4+c\sin\theta}
 &=\int_0^1\frac{dt}{2t^2+ct+2}\\
 &=\frac1d
   \left[\log\!\left(\frac{4t+c-d}{4t+c+d}\right)\right]_{t=0}^{t=1}\\
 &=3\sqrt2\log\!\left(\frac3{2\sqrt2}\right).
\end{align*}
Multiplication by $1/(6\sqrt6)$ is exactly the constant in
\eqref{koide-param-eq:I-infty}.
If the strict chamber wall fails, then the calculation in the proof of
Theorem~\ref{koide-param-thm:integral-parametrization}, now over $\RR$, gives
$q^2+4qn+n^2\leq1$.  The removed real subset is therefore compact.
For all large $X$ it contributes a fixed finite measure, proving
\eqref{koide-param-eq:volume-linear-O1}.

It remains to check effective well-roundedness for the actual sector.
Let $U_\delta$ be the image in the fixed vector representation of the
group ball used by Benoist--Oh~\cite[Definition~13.1]{BenoistOh_2012}.
For every linear functional $\ell$ on
$\RR^3$ there is a constant, independent of $x$ and $\delta$, such that
\begin{equation}
 |\ell(xg)-\ell(x)|\leq C_\ell\delta\|x\|
 \qquad(g\in U_\delta).
 \label{koide-param-eq:linear-functional-motion}
\end{equation}
Apply this to the five literal boundary functionals
\[
 a,\qquad q,\qquad n,\qquad a+4n-q,\qquad4a+17n.
\]
On the unbounded part of the positive sheet, $\|x\|\asymp a+1$, and the
last functional is also comparable to $\|x\|$.  Thus membership in the
outer height inequality can change only in a shell of radial width
$O(\delta X)$.  Membership in $q\geq0$ or $n\geq0$ can change only in
angular strips of width $O(\delta)$ about $\theta=\pi/2$ or $\theta=0$.
The boundary $a=0$ is compact.  The boundary $a+4n-q=0$ is compact by
the calculation preceding \eqref{koide-param-eq:volume-linear-O1}; moreover its
restriction to $\mathscr X$ is regular, because proportionality of
$(1,-1,4)$ to $(2a,-6q,-36n)$ is incompatible with $a,q,n\geq0$.
Consequently its group-action boundary tube has measure $O(\delta)$.
By \eqref{koide-param-eq:polar-measure}, all these regions together have measure
$O(\delta X)$.  The same analysis applies to the inner intersection that
defines $B^-_{X,\delta}$.  Hence
\begin{equation}
 m_\infty(B^+_{X,\delta}\setminus B^-_{X,\delta})
 \ll \delta\,m_\infty(B_X),
 \label{koide-param-eq:BO-boundary-condition}
\end{equation}
which is precisely Definition 13.1(2), with boundary exponent $1$, for
the group-action thickenings rather than merely Euclidean tubes.

For Definition~13.1(3) of \cite{BenoistOh_2012}, take the archimedean height to be the
Euclidean norm in the displayed integral basis.  It is bounded below on
$\mathscr X$.  Uniformly for $0<\delta<1$, the thickening
$B^+_{X,\delta}$ lies in a fixed finite enlargement of the polar sector;
there $H_\infty(x)\asymp1+a$ and the measure in a radial interval is
$O(da)$.  Hence for every $s>0$, uniformly in $\delta$,
\[
 \int_{B_X^+}\|u\|^{-s}\,dm_\infty(u)
 \ll
 \begin{cases}
 X^{1-s},&0<s<1,\\
 \log X,&s=1,\\
 1,&s>1.
 \end{cases}
\]
Each line is $O(m_\infty(B_X)^{1-\delta_s})$, for example with
$\delta_s=\min(s/2,1/2)>0$.  This is the remaining condition.  Since
$S=\{\infty\}$, Definition 13.1(1), concerning a finite-place compact
open subgroup, is void.  The choices of strict
or weak inequality on the finitely many piecewise smooth boundary arcs do
not change these neighborhood estimates.
\end{proof}

\begin{theorem}[Unconditional effective conclusion]
\label{koide-param-thm:power-saving}
There exists $\eta>0$, depending only on the fixed form, lattice, chamber,
and height in \eqref{koide-param-eq:original-problem}, such that
\begin{equation}
 R(X)=C_0X+O(X^{1-\eta}).
 \label{koide-param-eq:power-saving-again}
\end{equation}
\end{theorem}

\begin{proof}
Use the standard vector representation
$\rho:\operatorname{Spin}(F)\to\operatorname{SO}(F)\subset
\operatorname{GL}_3$, whose kernel is finite, and the displayed integral
basis of the exact lattice $\ZZ^3$.  Let
\[
 \mathbf G_{\ZZ}
 =\{g\in\mathbf G(\QQ):\rho(g)\ZZ^3=\ZZ^3\}.
\]
This is the integral group attached to the chosen representation and
lattice in the Benoist--Oh setup~\cite[pp.~1892--1893]{BenoistOh_2012};
take their congruence subgroup
$\Gamma_S$ to be the level-one group $\mathbf G_{\ZZ}$ itself.  The group
$\mathbf G$ is connected,
simply connected, and almost $\QQ$-simple.  It is noncompact over $\RR$
because $F$ has signature $(1,2)$.  Thus $S=\{\infty\}$ contains every
noncompact archimedean place and $G_S$ is noncompact.  The orbit
$x_0\mathbf G$ is the affine symmetric variety $\mathscr X$, and its
stabilizer has no nontrivial rational character by
\eqref{koide-param-eq:stabilizer-complement}.  These verify every standing hypothesis
on pp.~1892--1893 of \cite{BenoistOh_2012}.  Their Theorem~12.5 supplies the required
effective equidistribution of the symmetric-stabilizer translates.

The Borel--Harish-Chandra finiteness theorem quoted in \cite[p.~1892]{BenoistOh_2012}
p.~1892 says that $\mathscr X(\ZZ)$ is a finite union of
$\mathbf G_{\ZZ}$-orbits.  A rational representative of any such orbit
again has a symmetric stabilizer with no rational
character: its orthogonal complement is anisotropic over $\QQ$, since an
isotropic vector in it would be an isotropic vector for $F$.  Apply
Benoist--Oh, Theorem~13.2~\cite[pp.~1931--1932]{BenoistOh_2012}, to each orbit and to the exact effectively
well-rounded family $B_X$ of Lemma~\ref{koide-param-lem:volume}.  For each orbit it
gives its invariant-volume main term with error
$O(m_\infty(B_X)^{1-\eta_j})$ for some $\eta_j>0$.  There are finitely
many orbits, so their sum has the same form with
$\eta=\min_j\eta_j>0$.

All invariant measures on the same real homogeneous variety are
proportional.  Hence the sum of the finitely many orbit main terms is
$K\,m_\infty(B_X)$ for one fixed $K>0$.  Lemma~\ref{koide-param-lem:volume} and the
already established asymptotic $R(X)\sim C_0X$ force
$K I_\infty=C_0$.  Therefore
\[
 K\,m_\infty(B_X)=K(I_\infty X+O(1))=C_0X+O(1).
\]
Since $m_\infty(B_X)\asymp X$, summing the effective errors proves
\eqref{koide-param-eq:power-saving-again}; the $O(1)$ term is absorbed after decreasing
$\eta$ if necessary.  The theorem is applied first to the sequence of
positive integral $X$.  For real $X$, the original height $4a+17n$ is an
integer, so $R(X)=R(\lfloor X\rfloor)$; replacing $X$ by $\lfloor X\rfloor$
changes $C_0X$ by $O(1)$ and proves the same estimate.
\end{proof}

\begin{remark}[What exponent has actually been obtained]
Theorem~\ref{koide-param-thm:power-saving} is an unconditional exponent strictly
smaller than $1$: the error is $O(X^{1-\eta})$ for some $\eta>0$.
Neither the statement nor the proof of \cite[Theorem~13.2]{BenoistOh_2012} identifies an exponent
$1-\eta\leq1/2$ for this sharp sector.  Tracing their Sobolev and mixing
constants would make a positive $\eta$ effective in principle, but that
calculation has not been carried out in this section and cannot be replaced
by the assertion $\eta=1/2$.
\end{remark}

\section{Exact obstruction to the square-root claim}

There are two complementary descriptions of the remaining obstruction;
the first is an exact identity, while the second explains the relevant
spectral mechanisms without asserting an unproved intertwining formula.

First, the elementary ideal description gives the literal identity
\eqref{koide-param-eq:exact-discrepancy-splitting}.  Thus the square-root claim is
equivalent to a bound of that size for the displayed sum of the radial
bracket, the aggregate phase discrepancy, and $R_0(X)$.  Controlling the
radial and phase terms separately is a sufficient route, though the exact
identity also permits cancellation between them.  Formula
\eqref{koide-param-eq:Fejer-exact} shows that a separate estimate of the latter entails
uniform cancellation in the nonzero archimedean Hecke modes
\eqref{koide-param-eq:polynomial-Hecke-samples}.  Merely using
\(\#\mathcal P(m_n)=\rho(m_n)\), or replacing each phase distribution by
Lebesgue measure, deletes a term that is exactly present in $R(X)$.

Second, $F$ is rationally anisotropic, so $\mathbf G$ gives a cocompact
arithmetic Fuchsian quotient.  A negative-norm vector represents an
oriented geodesic, and its stabilizer is the split real torus appearing in
\eqref{koide-param-eq:stabilizer-complement}.  In the standard radial torus-orbit
counting kernel, the spectral expansion has two possible sources larger
than the square-root scale:
\begin{enumerate}
 \item an exceptional eigenvalue $\lambda=s(1-s)<1/4$ can contribute
 $c_sX^s$ with $s>1/2$, unless the exact torus period and exact sector
 transform make $c_s=0$;
 \item after exceptional terms are absent or removed, estimating the
 tempered spectral sum for a discontinuous cutoff at the conjectural
 $X^{1/2+\e}$ scale requires cancellation beyond the general
 effective-mixing argument.
\end{enumerate}
For the present nonradial chamber, one must first construct a
remainder-preserving sector spectral expansion before these candidate
terms can even be asserted with their exact coefficients.  An actual
refutation by a secondary term would then require proving both $s>1/2$
and $c_s\ne0$ for this fixed arithmetic quotient and exact chamber.  No
such construction or nonvanishing has been established here.  Conversely,
absence of exceptional terms alone would not estimate the full tempered
sharp-cutoff sum.  This is why the exact parametrization neither proves
nor refutes \eqref{koide-param-eq:square-root-target}, while it does isolate the data a
future proof must bound.

For comparison only, not as a theorem about the present nonradial sector,
the classical hyperbolic lattice-point problem for a conjugacy class has
an unconditional $O(X^{2/3})$ remainder due to Good
\cite[Theorem~4]{Good_1983}, with the modern proof of
Chatzakos--Petridis~\cite[Theorem~1.1 and Conjecture~5.7]{ChatzakosPetridis_2016},
and a conjectural $O_\e(X^{1/2+\e})$ scale.  That
published $O(X^{2/3})$ theorem counts a different radial kernel.  An exact
intertwining of its kernel with the simultaneous chamber and height in
\eqref{koide-param-eq:original-problem} has not been proved, so $O(X^{2/3})$ is not
claimed here for $R(X)$.

The independent finite replay used to check the HNF multiplicities, seed
reduction, inverse descent, and divisor factorization for \(1\leq n\leq200\)
is retained at
\path{supporting_materials/workbench/evidence/koide_exact_checkpoint_20260904/parametrization/verify_parametrization.py};
its SHA--256 digest is
\texttt{061c5e22258628169fd1ecb8bcc5e7cf\allowbreak%
8b8d19bfdb918650668c3e90f2f9730d}.
That replay is a finite consistency check.  The proofs above, rather than the
finite range, establish the parametrization for every admissible integer.

\section{Boundary ledger}

For clarity, every boundary used above has the following exact treatment.
\begin{itemize}
 \item $a>0$ is strict.  There is no omitted integral solution with
 $a=0$, because the original equation would give a multiple of $3$ equal
 to $2$.
 \item $q=0$ is included.  It corresponds to phase $\tau=0$, is included
 by $\tau\leq\lambda$, and is retained in the endpoint term
 $b_n(0)/2$ in \eqref{koide-param-eq:Fejer-exact}.
 \item $n=0$ is included and is counted exactly by $R_0(X)$; the sole
 solution $(a,q,n)=(1,1,0)$ on $a+4n=q$ is excluded.
 \item $a+4n=q$ is strict.  Its real intersection is bounded and its sole
 integral point in the nonnegative chamber is $(1,1,0)$.
 \item $4a+17n=X$ is included.  Equality is included in
 \eqref{koide-param-eq:t-height-bound}, \eqref{koide-param-eq:k-phase-bound}, and every floor in
 \eqref{koide-param-eq:exact-floor-formula}.
 \item The canonical unit interval is half-open,
 $[\sqrt{m_n},\epsilon\sqrt{m_n})$.  Consequently no generator is counted
 twice at a unit boundary.
 \item Conjugate ideals are distinct unless the ideal is fixed by
 conjugation.  They yield distinct canonical phases and distinct
 $q\geq0$ unit-orbits; no division by $2$ is made.  When $q=0$, the ideal
 is self-conjugate and still occurs once by HNF uniqueness.
\end{itemize}
